\subsection{Regularization and Model Selection}

Regularization controls model complexity by penalizing unstable parameter directions or by constraining the feasible set.

\begin{definition} \label{def:sl-penalized-linear-estimator} \hlt{Penalized Linear Estimation}\\
For empirical loss $L_n(\beta)$, a penalized estimator solves
\[
\widehat{\beta}_{\lambda}\in
\operatorname*{arg\,min}_{\beta}
\{L_n(\beta)+\lambda\Omega(\beta)\}.
\]
The choices $\Omega(\beta)=\norm{\beta}_2^2$ and $\Omega(\beta)=\norm{\beta}_1$ give ridge and lasso penalties.
\end{definition}

\begin{remark} \label{rem:sl-ridge-lasso-geometry} \hlt{Penalty Geometry}\\
The $\ell^2$ penalty shrinks coefficients smoothly along directions of weak design information.  The $\ell^1$ penalty has corners and can set coefficients exactly to zero, linking prediction to variable selection.
\end{remark}

\begin{definition} \label{def:sl-model-selection-validation} \hlt{Model Selection by Validation}\\
A validation criterion estimates predictive risk for candidate models or hyperparameters.  In $K$-fold cross-validation, the data are split into $K$ folds; each fold is held out once while the model is trained on the remaining folds, and the held-out losses are averaged.
\end{definition}

\begin{remark} \label{rem:sl-selection-bias} \hlt{Selection Bias}\\
The same validation data should not be used repeatedly as if it were untouched test data.  Nested validation or a final test set is needed when model selection itself is part of the learning procedure.
\end{remark}
